\section{Практичне використання криптографічних бібліотек}
\subsection{Симетричне автентифіковане шифрування}

Встановлення бібліотеки.
\begin{lstlisting}[language=Python]
import os
from cryptography.hazmat.primitives.ciphers.aead import AESGCM
from cryptography.exceptions import InvalidTag
\end{lstlisting}

\subsubsection{Генерація криптографічно стійкого ключа}

\begin{lstlisting}[language=Python]
key = AESGCM.generate_key(bit_length=256)
aesgcm = AESGCM(key)

message = b"Hello, my name is Mariia"
\end{lstlisting}

\subsubsection{Шифрування тестового повідомлення}

\begin{lstlisting}[language=Python]
nonce = os.urandom(12)
ciphertext = aesgcm.encrypt(nonce, message, None)

print("Message:", message)
print("Key:", key.hex())
print("Nonce:", nonce.hex())
print("Ciphertext:", ciphertext.hex())

\end{lstlisting}

\subsubsection{Розшифрування повідомлення}

\begin{lstlisting}[language=Python]
decrypted = aesgcm.decrypt(nonce, ciphertext, None)

print("\nDecrypted message:", decrypted)
\end{lstlisting}

\subsubsection{Перевірка цілісності після зміни шифротексту}

\begin{lstlisting}[language=Python]
modified_ciphertext = bytearray(ciphertext)
modified_ciphertext[0] ^= 1
modified_ciphertext = bytes(modified_ciphertext)

print("\nModified ciphertext:", modified_ciphertext.hex())

try:
    aesgcm.decrypt(nonce, modified_ciphertext, None)
    print("Decryption successful")
except InvalidTag:
    print("Error: authentication tag verification failed")
\end{lstlisting}

\subsubsection{Вивід}
У результаті виконання програми було згенеровано 256-бітний ключ, 12-байтовий nonce та отримано шифротекст тестового повідомлення.

\begin{lstlisting}[language=Python]
Message: b'Hello, my name is Mariia'
Key: fc5febdc3df5cd09ecb790ebea5651e7dd441783d988c1e643e58acc8db8f25d
Nonce: 9f72c8ca0d7e717102b1c728
Ciphertext: 1e256814693b724c81fb72b20cb1ee6bcae3db883a2a5fb36416e6af01e720ed626effe2d343afd2

Decrypted message: b'Hello, my name is Mariia'

Modified ciphertext: 1f256814693b724c81fb72b20cb1ee6bcae3db883a2a5fb36416e6af01e720ed626effe2d343afd2
Error: authentication tag verification failed
\end{lstlisting}

Після зміни одного байта шифротексту його розшифрування завершилося помилкою InvalidTag. Це свідчить про виявлення порушення цілісності зашифрованих даних.

\subsubsection{Призначення ключа, nonce та authentication tag}

У використаному алгоритмі AES-GCM ключ є секретним значенням, за допомогою якого виконується шифрування 
та розшифрування повідомлення. У цьому прикладі використовується ключ довжиною 256 біт.

Nonce --- це значення, яке використовується разом із ключем під час шифрування. 
У нашому випадку nonce має довжину 12 байт (Для AES-GCM рекомендована довжина nonce становить 12 байт (96 біт))
і генерується за допомогою криптографічно стійкого генератора випадкових чисел.
Для AES-GCM важливо, щоб nonce не повторювався при використанні одного й того самого ключа. 
Повторне використання nonce з тим самим ключем є небезпечним, оскільки може призвести до розкриття інформації про відкритий текст та порушення автентифікаційного захисту. 
Тому для кожного нового повідомлення необхідно використовувати новий nonce.

Authentication tag --- це тег, який генерується під час шифрування та використовується для 
перевірки цілісності й автентичності шифротексту. У використаній бібліотеці він формується автоматично.

Після зміни одного байта шифротексту перевірка authentication tag не проходить, тому виникає  \texttt{InvalidTag}. 
Це показує,що AES-GCM забезпечує не тільки конфіденційність даних, а й перевірку їх цілісності.

\subsection{Цифровий підпис}

Для практичної демонстрації створення та перевірки цифрового підпису використано алгоритм на основі скрученої кривої Едвардса --- Ed25519 (схема EdDSA). Він забезпечує високу швидкодію, стійкість до атак за часом виконання та фіксовану довжину підпису (64 байти) і відкритого ключа (32 байти).

Імпорт необхідних модулів:
\begin{lstlisting}[language=Python]
from cryptography.hazmat.primitives.asymmetric import ed25519
from cryptography.hazmat.primitives import serialization
from cryptography.exceptions import InvalidSignature
\end{lstlisting}

\subsubsection{Генерація ключової пари}

\begin{lstlisting}[language=Python]
private_key = ed25519.Ed25519PrivateKey.generate()
public_key = private_key.public_key()

pub_bytes = public_key.public_bytes(
    encoding=serialization.Encoding.Raw,
    format=serialization.PublicFormat.Raw
)

message = b"Hello, my name is Mariia"

print("Private key: generated (held in memory, not exported)")
print("Public key (raw hex):", pub_bytes.hex())
print("Original message:", message)
\end{lstlisting}

\subsubsection{Створення цифрового підпису}

\begin{lstlisting}[language=Python]
signature = private_key.sign(message)

print("\nSignature (hex):", signature.hex())
print("Signature length:", len(signature), "bytes")
\end{lstlisting}

\subsubsection{Перевірка цифрового підпису}

\begin{lstlisting}[language=Python]
try:
    public_key.verify(signature, message)
    print("Verification result: Signature is VALID")
except InvalidSignature:
    print("Verification result: Signature is INVALID")
\end{lstlisting}

\subsubsection{Перевірка підпису після зміни повідомлення}

\begin{lstlisting}[language=Python]
modified_message = b"Hello, my name is Mariia!"

print("\nModified message:", modified_message)

try:
    public_key.verify(signature, modified_message)
    print("Verification result: Signature is VALID")
except InvalidSignature:
    print("Verification result: Signature verification FAILED (InvalidSignature)")
\end{lstlisting}

\subsubsection{Вивід}

У результаті виконання коду було згенеровано пару ключів Ed25519, обчислено цифровий підпис вихідного повідомлення та успішно верифіковано його за допомогою відкритого ключа.

\begin{lstlisting}[language=bash]
Private key: generated (held in memory, not exported)
Public key (raw hex): 25aed4306658c129d87a34fff7b6497f09f3c8e30a5f694f69912c5b1ff45b88
Original message: b'Hello, my name is Mariia'

Signature (hex): 9b304a2c75ea4b89fffab3bd995d40930b295f1d012d4b279aaab3906b560ea6a3edc889198a5a15
    454d2bd50b86103e060d23edafcda683ad5e48d3ccb44504
Signature length: 64 bytes
Verification result: Signature is VALID

Modified message: b'Hello, my name is Mariia!'
Verification result: Signature verification FAILED (InvalidSignature)
\end{lstlisting}

\subsubsection{Призначення ключів та цифрового підпису}

У криптосистемі Ed25519 закритий ключ є секретним і застосовується виключно для накладання цифрового підпису, тоді як відкритий ключ поширюється вільно та дозволяє будь-якій стороні перевірити справжність документа. Цифровий підпис фіксованого розміру (64 байти) підтверджує авторство відправника та незмінність даних. Модифікація навіть одного знака у повідомленні (додавання знака оклику) призвела до завершення перевірки винятком \texttt{InvalidSignature}, що підтверджує надійне виявлення спотворення підписаних даних.
